\documentclass[11pt]{article}
% Source-PDF: IOI中国国家候选队论文/2004/金恺.pdf
\usepackage{oipl-vn}

\newcommand{\Oh}{\mathcal{O}}
\newcommand{\area}[1]{S_{#1}}
\newcommand{\dist}{\operatorname{dist}}
\newcommand{\abs}[1]{\left|#1\right|}

\title{Phương pháp giới hạn: con đường ngắn để giải các bài toán tối ưu hình học}
\author{Jin Kai\\Trường Trung học Changjun}
\date{IOI 2004, đội tuyển tập huấn quốc gia Trung Quốc}

\begin{document}
\maketitle

\origtitle{The limit method: a shortcut for geometric optimization problems}
\sourcepath{IOI China candidate team papers / 2004 / Jin Kai.pdf}

\section*{Từ khóa}

Giá trị tối ưu; giới hạn; vô cùng bé; hiệu chỉnh vi mô.

\section*{Tóm tắt}

Trong các bài toán hình học phẳng, ta thường gặp yêu cầu tìm cực trị. Ở những
bài toán đó, biến và hàm mục tiêu có thể liên quan đến tọa độ, hệ số góc, độ
dài, góc, chu vi, diện tích, cùng nhiều đại lượng phức tạp khác. Bản thân biến
cũng thường chịu các ràng buộc hình học khó xử lý, nên cách trực tiếp lập hàm
rồi tìm cực trị hoặc không khả thi, hoặc cực kỳ rườm rà.

Trong nhiều bài toán như vậy, biến có vô số cách nhận giá trị; chẳng hạn một
điểm có thể chạy trên cả mặt phẳng hoặc trên một đường thẳng. Vì thế ta không
thể liệt kê mọi phương án để tìm giá trị tốt nhất. Khi đó thường có thể dùng
\term{phương pháp giới hạn} để chứng minh rằng: nếu biến nhận một giá trị không
thuộc một số trường hợp đặc biệt, thì hàm mục tiêu không thể tối ưu, bởi chỉ
cần hiệu chỉnh một lượng vô cùng bé\footnote{Ví dụ quay đi một góc vô cùng bé
hoặc tịnh tiến một đoạn vô cùng bé. Trong bài này, ``lượng vi mô'' luôn chỉ
một lượng vô cùng bé: miễn thay đổi đủ nhỏ, quan hệ tương đối giữa các điểm,
tình trạng giao nhau của các đoạn thẳng, v.v. sẽ không đổi.} đã có thể làm
giá trị hàm mục tiêu tốt hơn. Như vậy chỉ còn hữu hạn trường hợp đặc biệt có
thể là nghiệm tối ưu; liệt kê chúng là đủ để tìm nghiệm tối ưu.\footnote{Về
bản chất, phương pháp giới hạn chính là xét đạo hàm của hàm mục tiêu. Nếu đạo
hàm khác \(0\) và biến không nằm trên biên miền xác định, điểm đó không thể là
cực trị. Đây cũng là lý do dùng tên ``phương pháp giới hạn''.}

Trong một số bài toán cùng loại khác, về nguyên tắc đã có thể liệt kê hữu hạn
trường hợp để tìm tối ưu, nhưng lượng cần liệt kê rất lớn và độ phức tạp cao.
Ta cũng có thể thử dùng phương pháp giới hạn để giảm mạnh số trường hợp cần
xét, từ đó hạ độ phức tạp.

Tóm lại, ý tưởng chung của phương pháp giới hạn là biến vô hạn thành hữu hạn,
và biến hữu hạn lớn thành một lượng nhỏ.

\section{Mở đầu}

Phương pháp giới hạn có rất nhiều ví dụ ứng dụng. Chẳng hạn trong bài toán
kinh điển \term{hình chữ nhật phủ nhỏ nhất}\footnote{Cho \(n\) điểm trên mặt
phẳng, tìm một hình chữ nhật diện tích nhỏ nhất sao cho mọi điểm đã cho nằm
bên trong nó.}, ta dùng phương pháp giới hạn để chứng minh rằng một cạnh của
hình chữ nhật nhỏ nhất phải đi qua hai điểm đã cho; nhờ vậy số hình chữ nhật
cần liệt kê giảm đi rất nhiều.

Dù vậy, phương pháp giới hạn thật sự không dễ dùng. Có lúc chứng minh rất khó,
các trường hợp có thể phức tạp, hoặc ta không biết phải hiệu chỉnh như thế nào.
Đôi khi còn phải dùng lượng giác và các kiến thức toán tương đối nặng. Vì thế,
muốn nắm được phương pháp này cần nền tảng toán vững, khả năng quan sát tốt và
một lượng kinh nghiệm nhất định.

Sau đây ta dùng phương pháp giới hạn để giải một số bài toán điển hình, qua đó
từng bước hiểu sâu hơn công dụng và cách sử dụng của nó.

\section{Ví dụ 1: Chocolate}

\subsection*{Mô tả bài toán}

Một nhà máy kẹo có loại chocolate dạng đa giác lồi \(N\) cạnh
\((4\le N\le 50)\). Kiddy và Carlson góp tiền mua một miếng, muốn cắt nó thành
hai nửa bằng đúng một nhát cắt. Hai phần phải có diện tích bằng nhau. Hãy tìm
độ dài ngắn nhất của đoạn cắt dùng để chia chocolate.

\subsection*{Mô hình toán học}

Cho \(N\) điểm \((X_i,Y_i)\), \(1\le i\le N\), tạo thành một bao lồi \(P\).
Cần tìm một đường cắt \(L\) sao cho diện tích hai phía của đường cắt bằng nhau,
và độ dài của \(L\) là nhỏ nhất.

\subsection*{Phân tích}

Gọi hai đầu mút của đoạn cắt là \(A,B\). Hai điểm này có thể là đỉnh của \(P\),
cũng có thể chỉ nằm trong phần trong của một cạnh của \(P\). Ta tách hai loại
trường hợp.

\begin{enumerate}
  \item \(A\) là đỉnh của \(P\) (hoặc \(B\) là đỉnh, khi đó đổi vai trò \(A,B\)).
  Ta hiển nhiên có thể liệt kê một đỉnh của \(P\) làm \(A\). Từ ràng buộc diện
  tích hai phía bằng nhau, vị trí của \(B\) được xác định trực tiếp.

  \item \(A,B\) đều nằm trên các cạnh của \(P\). Liệt kê hai cạnh chứa \(A\) và
  \(B\), gọi giao điểm của hai đường thẳng chứa chúng là \(C\).
\end{enumerate}

Đặt \(\angle C=\gamma\), \(\angle CAB=\alpha\),
\(\angle CBA=\beta\). Khi \(\alpha\ne\beta\), ta chứng minh được rằng nếu quay
đoạn cắt \(L\) một góc vô cùng bé \(\theta\) thành \(L'\), đồng thời vẫn giữ
điều kiện diện tích hai phía bằng nhau, thì có thể làm \(L'\) ngắn hơn \(L\).

Chọn hướng quay sao cho \(\pi>\beta>\alpha>0\). Sau khi quay, \(L'\) vẫn cắt
hai cạnh ban đầu, vì ta chỉ quay một lượng vô cùng bé. Gọi hai giao điểm mới là
\(A',B'\); khi đó
\[
  \angle CA'B'=\alpha+\theta,\qquad
  \angle CB'A'=\beta-\theta.
\]
Theo định lý sin trong tam giác \(ABC\),
\[
  \frac{AC}{\sin\beta}
  =\frac{BC}{\sin\alpha}
  =\frac{L}{\sin\gamma}.
\]
Trong tam giác \(A'B'C\),
\[
  \frac{A'C}{\sin(\beta-\theta)}
  =\frac{B'C}{\sin(\alpha+\theta)}
  =\frac{L'}{\sin\gamma}.
\]
Vì \(L\) và \(L'\) đều là đường chia đôi diện tích, ta có
\(\area{\triangle ABC}=\area{\triangle A'B'C}\), tức
\[
\begin{aligned}
  \frac12 AC\cdot BC\sin\gamma
    &=\frac12 A'C\cdot B'C\sin\gamma,\\
  AC\cdot BC&=A'C\cdot B'C.
\end{aligned}
\]
Thay các hệ thức định lý sin vào:
\[
\begin{aligned}
  \frac{L\sin\beta}{\sin\gamma}\cdot
  \frac{L\sin\alpha}{\sin\gamma}
  &=
  \frac{L'\sin(\beta-\theta)}{\sin\gamma}\cdot
  \frac{L'\sin(\alpha+\theta)}{\sin\gamma},\\
  \frac{L^2}{L'^2}
  &=
  \frac{\sin(\beta-\theta)\sin(\alpha+\theta)}
       {\sin\beta\sin\alpha}\\
  &=
  \frac{\cos(\beta-\alpha-2\theta)-\cos(\beta+\alpha)}
       {\cos(\beta-\alpha)-\cos(\beta+\alpha)}.
\end{aligned}
\]
Với \(\theta\) dương đủ nhỏ, từ \(\pi>\beta>\alpha>0\) suy ra
\[
  \cos(\beta-\alpha-2\theta)>\cos(\beta-\alpha),
\]
nên
\[
  \frac{\cos(\beta-\alpha-2\theta)-\cos(\beta+\alpha)}
       {\cos(\beta-\alpha)-\cos(\beta+\alpha)}>1.
\]
Do đó \(L^2>L'^2\), tức \(L'<L\). Nếu hai cạnh chứa \(A,B\) song song, có thể
xem \(C\) ở vô cực và lập luận tương tự vẫn đúng.

Vì vậy, nếu hai góc mà \(L\) tạo với hai cạnh chứa đầu mút không bằng nhau, thì
\(L\) không thể là đoạn cắt ngắn nhất. Nếu \(L\) là đoạn cắt ngắn nhất mà không
đi qua đỉnh nào của \(P\), thì hai góc đó bắt buộc bằng nhau.

Đây chính là kết luận ta cần. Sau khi liệt kê hai cạnh chứa hai đầu mút của
\(L\), hệ số góc của \(L\) đã được xác định; từ điều kiện diện tích hai phía
bằng nhau, ta tính trực tiếp vị trí của \(L\). Nhờ các kết luận trên, ta có thể
thiết kế thuật toán \(\Oh(N^2)\).

\paragraph{Tiểu kết.}
Qua bài này, ta lần đầu tiếp xúc với phương pháp giới hạn và dùng nó thu được
một kết luận rất đơn giản. Miền giá trị của biến được giảm từ vô hạn đường
thẳng xuống hữu hạn trường hợp, sau đó bài toán được giải bằng liệt kê đơn
giản.

\section{Ví dụ 2: Space (bài thảo luận tập huấn 2003 số 0039)}

\subsection*{Mô tả bài toán}

Trên mặt phẳng có \(n\) điểm đôi một khác nhau \((3\le n\le 10000)\). Cần tìm
một đường thẳng sao cho tổng khoảng cách từ mọi điểm đến đường thẳng đó là nhỏ
nhất.

\subsection*{Mô hình toán học}

Cho tọa độ \(n\) điểm
\[
  V_1(x_1,y_1),V_2(x_2,y_2),\ldots,V_n(x_n,y_n).
\]
Với đường thẳng \(l: ax+by+c=0\), định nghĩa
\[
  f(l)=\sum_{i=1}^n
  \frac{\abs{ax_i+by_i+c}}{\sqrt{a^2+b^2}}.
\]
Mục tiêu là tìm \(\min f(l)\). Bản trích xuất PDF ghi điều kiện \(ab\ne0\);
về ý nghĩa hình học, công thức trên chỉ cần \((a,b)\ne(0,0)\).

\subsection*{Phân tích}

Cách nghĩ dễ nhất là liệt kê mọi đường thẳng rồi lấy giá trị tốt nhất. Nhưng
trên mặt phẳng có vô hạn đường thẳng. Vậy đường thẳng nào mới có khả năng là
đường ta cần?

\paragraph{Kết luận 1.}
Có thể yêu cầu \(l\) đi qua một điểm đã cho.

Thật vậy, nếu \(l\) không đi qua điểm đã cho nào, một trong hai phía của \(l\)
có số điểm không ít hơn phía còn lại. Gọi số điểm phía nhiều hơn là \(a\), phía
ít hơn là \(b\). Tịnh tiến \(l\) về phía có nhiều điểm hơn một lượng vi mô
\(\Delta\), thu được \(l'\). Khi đó
\[
  f(l')-f(l)=b\Delta-a\Delta=(b-a)\Delta\le0,
\]
nên \(f(l')\le f(l)\). Cứ tịnh tiến theo cùng hướng cho đến khi gặp điểm đã
cho đầu tiên, ta thu được \(l''\) với \(f(l'')\le f(l)\).

\paragraph{Kết luận 2.}
Có thể tiếp tục yêu cầu \(l\) đi qua hai điểm đã cho.

Theo kết luận 1, giả sử \(l\) đi qua \(V_1\) nhưng không đi qua điểm nào khác.
Ký hiệu \(L_i\) là khoảng cách từ \(V_i\) đến \(V_1\), và \(\alpha_i\) là góc
giữa đoạn \(V_iV_1\) và đường vuông góc từ \(V_i\) xuống \(l\).

Quay \(l\) quanh \(V_1\) ngược chiều kim đồng hồ một góc rất nhỏ \(\alpha\)
\((\alpha\to0^+)\), được \(l'\); quay theo chiều kim đồng hồ cùng góc đó, được
\(l''\). Nếu \(\alpha\) đủ nhỏ, trong quá trình quay đường thẳng sẽ không gặp
điểm đã cho nào khác.

Nếu \(\alpha_i=0\), dù quay đến \(l'\) hay \(l''\), khoảng cách từ \(V_i\) đến
đường thẳng đều giảm nghiêm ngặt. Nếu \(\alpha_i\ne0\), hai góc sau khi quay
lần lượt là \(\alpha_i+\alpha\) và \(\alpha_i-\alpha\). Do
\[
  \cos(\alpha_i-\alpha)+\cos(\alpha_i+\alpha)
  =2\cos\alpha_i\cos\alpha<2\cos\alpha_i,
\]
ta có
\[
  L_i\cos(\alpha_i-\alpha)+L_i\cos(\alpha_i+\alpha)
  =2L_i\cos\alpha_i\cos\alpha<2L_i\cos\alpha_i.
\]
Cộng thay đổi của mọi điểm lại, suy ra
\[
  f(l')+f(l'')<2f(l).
\]
Nếu \(l\) là tối ưu, phải có \(f(l)\le f(l')\) và \(f(l)\le f(l'')\), nên
\[
  f(l')+f(l'')\ge f(l)+f(l),
\]
mâu thuẫn. Vì thế có thể quy định \(l\) phải đi qua hai điểm đã cho.

Đến đây, số đường thẳng cần liệt kê đã hữu hạn. Một thuật toán trực tiếp là:
\[
\begin{array}{l}
\text{min}\leftarrow\infty;\\
\text{liệt kê hai điểm;}\\
\quad\text{xác định đường thẳng }h\text{ qua hai điểm đó;}\\
\quad\text{tính }\text{now}\leftarrow f(h);\\
\quad\text{nếu }\text{now}<\text{min}\text{ thì cập nhật }l\leftarrow h,\ 
  \text{min}\leftarrow \text{now}.
\end{array}
\]
Phương pháp giới hạn đã biến tập đường thẳng cần xét từ vô hạn thành hữu hạn
\(\Theta(n^2)\) đường, nhờ đó ta có thuật toán hữu hiệu. Tuy nhiên thuật toán
này có độ phức tạp \(\Oh(n^3)\), còn cần giảm tiếp bằng cách loại các đường
thẳng vô dụng.

\paragraph{Kết luận 3.}
Định nghĩa \(a(l)\) là số điểm phía trên \(l\) (nếu \(l\) thẳng đứng thì là
phía bên phải), \(b(l)\) là số điểm nằm trên \(l\), và \(c(l)\) là số điểm
phía dưới \(l\) (nếu \(l\) thẳng đứng thì là phía bên trái). Nếu \(l\) là tối
ưu, tất yếu
\[
  a(l)+b(l)>c(l),\qquad c(l)+b(l)>a(l).
\]

Chứng minh: nếu \(a(l)+b(l)\le c(l)\) hoặc \(c(l)+b(l)\le a(l)\), ta tịnh
tiến \(l\) một lượng vi mô về phía có nhiều điểm hơn (có thể bằng nhau), thu
được \(l'\). Theo chứng minh của kết luận 1, \(f(l')\le f(l)\). Vì chỉ di
chuyển một lượng vi mô, trên \(l'\) không có điểm đã cho nào. Lấy một điểm bất
kỳ \(A\) trên \(l'\), xem \(A\) như \(V_1\) trong kết luận 2, rồi quay vi mô
quanh \(A\). Bằng lập luận tương tự kết luận 2, \(f(l')\) không thể là tối ưu,
do đó \(f(l)\) cũng không thể là tối ưu.

Gọi \(E\) là tập các đường thẳng thỏa cả ba kết luận trên. Có thể chứng minh
\(\abs{E}\) cùng cấp với \(n\). Dùng phương pháp quay, từ một đường thẳng trong
\(E\) có thể tìm đường tiếp theo trong thời gian \(\Oh(n)\); quay một vòng sẽ
tìm được toàn bộ các đường trong \(E\). Việc tính \(f\) của một đường thẳng
cũng chỉ cần \(\Oh(n)\). Vì vậy độ phức tạp tổng cộng chỉ còn \(\Oh(n^2)\).%
\footnote{Chứng minh và thuật toán chi tiết được trình bày trong báo cáo bài
tập tập huấn năm 2003 của tác giả, bài 0039.}

\paragraph{Tiểu kết.}
Các chứng minh trong bài này không phức tạp. Về bản chất, cả ba kết luận đều
gắn chặt với phương pháp giới hạn:
\begin{itemize}
  \item Tịnh tiến một lượng vi mô để chứng minh một lớp đường thẳng không thể
  tối ưu, từ đó giảm miền giá trị của \(l\) từ mọi đường thẳng trên mặt phẳng
  xuống các đường thẳng đi qua một điểm đã cho.
  \item Quay một lượng vi mô để loại tiếp một lớp trong phần còn lại, giảm miền
  giá trị của \(l\) xuống hữu hạn đường.
  \item Tịnh tiến vi mô rồi quay vi mô để giảm số đường thẳng cần xét từ
  \(\Theta(n^2)\) xuống \(\Theta(n)\).
\end{itemize}

Từ bài này có thể thấy một quy luật thường gặp khi giải tối ưu hình học phẳng:
sau khi gặp bài toán, ta dễ nảy sinh phỏng đoán rằng nghiệm tối ưu có thỏa một
tính chất nào đó không; nếu có, liệu nó có thỏa một tính chất đặc biệt hơn nữa
không; cứ tiếp tục đặt và chứng minh phỏng đoán như vậy để thu hẹp miền giá
trị của biến, cho đến khi không thể thu hẹp thêm hoặc đã đủ tốt. Phần còn lại
chỉ cần liệt kê và tính toán.

Các phỏng đoán này có khi đúng, có khi có phản ví dụ; có phỏng đoán hiển nhiên,
có phỏng đoán chứng minh rất khó. Cách hình thành phỏng đoán đơn giản và hiệu
quả nhất là xét vài ví dụ nhỏ để tìm quy luật; trực giác và cảm hứng đến từ
quan sát cẩn thận. Cách chứng minh đầu tiên là kiểm tra trên vài ví dụ. Nếu có
phản ví dụ, phỏng đoán thất bại và ta phải sửa một phần hoặc đưa ra phỏng đoán
mới. Nếu không tìm được phản ví dụ, điều đó vẫn không có nghĩa phỏng đoán là
đúng; ta cần phân tích chặt chẽ để chứng minh đầy đủ. Phương pháp giới hạn
chính là một công cụ phân tích đơn giản và thực dụng.

Hai bài đầu khá đơn giản: ta trực tiếp đưa ra phỏng đoán đúng và chứng minh
cũng dễ. Nhưng khi giải bài thật, mọi việc thường không thuận lợi như vậy. Ta
sẽ xét trọn vẹn một bài phức tạp hơn.

\section{Ví dụ 3: Hệ thống phòng thủ tên lửa}

Đây là bài tự xây dựng của tác giả.

\subsection*{Mô tả bài toán}

Biên giới của một quốc gia là một đa giác lồi. Để phòng thủ, nhà vua dự định
cho xây 4 tháp phòng thủ. Khi chiến tranh nổ ra, kẻ địch có thể tập kích một
trong bốn tháp. Một khi tháp đó bị tập kích, ba tháp còn lại sẽ vào trạng thái
chiến đấu và phong tỏa hoàn toàn vùng tam giác tạo bởi ba đoạn nối từng cặp
trong số chúng. Kẻ địch xảo quyệt sẽ khiến diện tích vùng bị phong tỏa càng
nhỏ càng tốt, còn nhà vua thông minh muốn diện tích vùng còn hoạt động sau khi
một tháp bị phá càng lớn càng tốt.

\subsection*{Mô hình toán học}

Định nghĩa
\[
  f(A,B,C,D)=\min\{
  \area{\triangle ABC},\area{\triangle ABD},
  \area{\triangle ACD},\area{\triangle BCD}\}.
\]
Cho một đa giác lồi \(n\) cạnh \(P\) trên mặt phẳng. Hãy chọn 4 điểm
\(A,B,C,D\) trong \(P\), được phép nằm trên biên, sao cho
\(f(A,B,C,D)\) lớn nhất.

\subsection{Chuyển hóa bài toán}

Hàm \(f\) là min của 4 đại lượng, nên rất khó xử lý trực tiếp: tam giác nào
mới là nhỏ nhất? Xét bốn điểm tùy ý \(A,B,C,D\) trên mặt phẳng. Có ba kiểu
bố trí.

\begin{enumerate}
  \item Có ba điểm thẳng hàng. Khi đó \(f(A,B,C,D)=0\).

  \item Không có ba điểm thẳng hàng, nhưng một điểm nằm trong tam giác tạo bởi
  ba điểm còn lại. Giả sử \(D\) nằm trong tam giác \(ABC\). Khi đó
  \[
  \begin{aligned}
    f(A,B,C,D)
    &=\min\{\area{\triangle ABD},\area{\triangle ACD},
      \area{\triangle BCD},\\
    &\qquad\quad
      \area{\triangle ABD}+\area{\triangle ACD}
      +\area{\triangle BCD}\}\\
    &=\min\{\area{\triangle ABD},\area{\triangle ACD},
      \area{\triangle BCD}\}.
  \end{aligned}
  \]
  Do
  \[
    \area{\triangle ABD}+\area{\triangle ACD}
      +\area{\triangle BCD}=\area{\triangle ABC},
  \]
  ta có \(f(A,B,C,D)\le \area{\triangle ABC}/3\). Vì \(ABCD\) được chọn tùy
  ý, ta có thể lấy \(D\) là trọng tâm của \(\triangle ABC\), khi đó
  \(f(A,B,C,D)=\area{\triangle ABC}/3\). Vì vậy, trong loại ``một điểm nằm
  trong tam giác của ba điểm còn lại'', giá trị \(f\) lớn nhất bằng một phần
  ba diện tích tam giác lớn nhất nằm trong \(P\). Gọi đây là bài toán \((1)\).

  \item \(A,B,C,D\) là bốn đỉnh của một tứ giác lồi. Giả sử
  \(f(A,B,C,D)=\area{\triangle BCD}\). Qua \(D\) kẻ đường song song với \(BC\),
  qua \(B\) kẻ đường song song với \(DC\), hai đường cắt nhau tại \(A'\). Khi
  đó \(A'\) nằm trong hoặc trên biên tam giác \(ABD\).
\end{enumerate}

Chứng minh nhận xét cuối: vì
\[
  \area{\triangle BCA'}=\area{\triangle BCD}\le\area{\triangle BCA},
\]
nên \(A\) và \(B\) nằm ở hai phía của \(A'D\). Tương tự,
\[
  \area{\triangle A'CD}=\area{\triangle BCD}\le\area{\triangle ACD},
\]
nên \(A\) và \(D\) nằm ở hai phía của \(A'B\). Do đó \(A'\) nằm trong hoặc
trên biên \(\triangle ABD\), và vì thế cũng nằm trong \(P\).

Vì \(A'BCD\) là hình bình hành, giá trị \(f\) tại bốn đỉnh của nó bằng một nửa
diện tích hình bình hành. Hơn nữa
\[
  f(A',B,C,D)=\frac{\area{A'BCD}}2=\area{\triangle BCD}=f(A,B,C,D).
\]
Nói cách khác, có thể điều chỉnh \(ABCD\) thành bốn đỉnh của một hình bình
hành nằm trong \(P\) mà không làm đổi \(f\). Vì vậy, trong loại tứ giác lồi,
giá trị \(f\) lớn nhất bằng một nửa diện tích hình bình hành lớn nhất nằm trong
\(P\). Gọi đây là bài toán \((2)\).

Nhờ chuyển hóa này, bài toán max-min ban đầu trở thành bài toán cực đại một
đại lượng hình học cụ thể. Bài mới rõ ràng dễ xét hơn bài gốc.

\subsection{Bài toán tam giác lớn nhất}

Trước hết xét \((1)\), vì có vẻ đơn giản hơn, và thực tế cũng đơn giản hơn
nhiều. Giả sử \(A,B,C\) đã xác định, ta xem chúng cần thỏa điều kiện gì.

Ba điểm \(A,B,C\) đều có thể lấy là đỉnh của \(P\) mà không làm giảm diện tích;
nói cách khác, diện tích tam giác lớn nhất trong \(P\) bằng diện tích lớn nhất
của một tam giác có ba đỉnh là đỉnh của \(P\).

Chứng minh: qua \(A\), kẻ đường thẳng \(l\) song song với \(BC\). Trên \(l\),
hoặc ở phía của \(l\) xa \(BC\), chắc chắn tồn tại một đỉnh \(A'\) của \(P\).
Khoảng cách từ \(A'\) đến \(BC\) không nhỏ hơn khoảng cách từ \(A\) đến \(BC\),
nên
\[
  \area{\triangle A'BC}\ge \area{\triangle ABC}.
\]
Vì thế có thể dời \(A\) đến một đỉnh nào đó mà diện tích tăng hoặc không đổi.
Lặp lại với \(B,C\) là đủ.

Do đó ta có thể dùng thuật toán \(\Oh(N^3)\) liệt kê ba đỉnh của \(P\), tìm
diện tích tam giác lớn nhất \(S\); lời giải cho bài toán \((1)\) là \(S/3\).

\subsection{Bài toán hình bình hành lớn nhất}

Trường hợp tam giác dễ vì có thể khiến \(A,B,C\) đúng là ba đỉnh của \(P\).
Với hình bình hành thì cách tương tự rõ ràng sai: có thể không có bốn đỉnh nào
của \(P\) lập thành hình bình hành; nếu bắt \(A,B,C,D\) đều là đỉnh của \(P\)
thì có khi vô nghiệm.\footnote{Nếu không chuyển hóa thành hình bình hành mà
trực tiếp liệt kê bốn đỉnh của \(P\) làm \(A,B,C,D\) để tính \(f(A,B,C,D)\),
kết quả cũng không tối ưu. Ví dụ với
\[
N=6,\quad
P=\{(-9.60,0.95),(-7.12,3.15),(-0.74,4.72),
(7.33,0.98),(4.26,-2.41),(-6.37,-0.152)\},
\]
nghiệm tối ưu là \(26.35\), còn liệt kê bốn đỉnh chỉ cho giá trị lớn nhất
\(19.40\).}

Ta vẫn bắt đầu bằng cách thu hẹp miền giá trị của \(A,B,C,D\).

\paragraph{Kết luận 1.}
Bốn điểm \(A,B,C,D\) đều nằm trên biên của \(P\), không nhất thiết là đỉnh.

Chứng minh: nếu một điểm không nằm trên biên, giả sử là \(A\), hãy tịnh tiến
cạnh \(AD\) theo hướng \(DA\) một lượng vi mô\footnote{Nhỏ đến mức
\(A'D'\) vẫn nằm trong phần trong của \(P\), không kể biên.} thành \(A'D'\).
Khi đó
\[
  \area{A'BCD'}=\area{ABCD},
\]
và \(A',D'\) đều chưa nằm trên biên. Tiếp tục tịnh tiến cạnh \(A'D'\) theo
hướng \(BA'\) một lượng vi mô\footnote{Nhỏ đến mức \(A''D''\) vẫn nằm trong
phần trong của \(P\), không kể biên.} thành \(A''D''\), ta có
\[
  \area{A''BCD''}>\area{ABCD}.
\]
Vì vậy nếu \(A\) không nằm trên biên, \(ABCD\) không thể là hình bình hành diện
tích lớn nhất trong \(P\). Suy ra ở trạng thái tối ưu, cả bốn điểm đều nằm
trên biên \(P\).

\paragraph{Kết luận 2.}
Có thể lấy một trong \(A,B,C,D\) là đỉnh của \(P\) mà không làm giảm diện tích.

Trước khi chứng minh, ta cần hai khái niệm nhỏ. Cho hai đường thẳng không song
song \(l,g\) và một điểm \(E\) không nằm đồng thời trên hai đường. Nếu có điểm
\(B\in l\), điểm \(D\in g\), và đoạn \(BD\) nhận \(E\) làm tâm đối xứng, thì
gọi \(BD\) là \term{đường chia đôi của \((l,g,E)\)}.

\paragraph{Bổ đề 1.}
Đường chia đôi của \((l,g,E)\) là duy nhất.

Chứng minh: quay \(g\) quanh \(E\) một góc \(180^\circ\), được đường \(g'\).
Hai đường \(g\) và \(g'\) đối xứng tâm qua \(E\). Gọi \(B=l\cap g'\); vì
\(l\) cắt \(g\), nó cũng cắt đường song song \(g'\). Gọi \(D\) là điểm đối
xứng của \(B\) qua \(E\). Khi đó \(D\in g\), và \(BD\) là đường chia đôi của
\((l,g,E)\). Mặt khác, \(l\) và \(g'\) chỉ có một giao điểm, nên không thể có
đường chia đôi thứ hai.

\paragraph{Bổ đề 2.}
Giả sử \(E'\) và \(E''\) đối xứng nhau qua \(E\), và \(EE'\) là lượng vi mô.
Gọi \(BD,B'D',B''D''\) lần lượt là đường chia đôi của
\((l,g,E)\), \((l,g,E')\), \((l,g,E'')\). Khi đó \(B'\) và \(B''\) đối xứng
nhau qua \(B\); đồng thời \(D'\) và \(D''\) đối xứng nhau qua \(D\).

Chứng minh theo ba trường hợp:
\begin{enumerate}
  \item Nếu \(E'E''\parallel g\), thì \(B,B',B''\) trùng nhau. Từ tính chất
  đường trung bình dễ thấy
  \[
    DD'=2EE'=2EE''=DD'',
  \]
  mà \(D,D',D''\) cùng nằm trên \(g\), nên \(D'\) và \(D''\) đối xứng qua
  \(D\). Kết luận về \(B\) là hiển nhiên.

  \item Nếu \(E'E''\parallel l\), chứng minh tương tự.

  \item Nếu \(E'E''\) không song song với \(g\) và cũng không song song với
  \(l\), do \(E'\) và \(E''\) đối xứng qua \(E\), ta có \(E'E''=2EE''\). Đặt
  \(x=H(E'',g)\), trong đó \(H(V,l)\) là khoảng cách từ điểm \(V\) đến đường
  thẳng \(l\), và đặt \(y\) là thành phần của \(EE''\) theo phương vuông góc
  với \(g\). Khi đó
  \[
    H(B'',g)=2x,\qquad H(B,g)=2(x+y),\qquad H(B',g)=2(x+2y).
  \]
  Suy ra \(BB''=BB'\). Vì \(B,B',B''\) cùng nằm trên \(l\), nên \(B'\) và
  \(B''\) đối xứng qua \(B\). Xét tương tự khoảng cách theo phương vuông góc
  với \(l\), ta được \(DD''=DD'\), nên \(D'\) và \(D''\) đối xứng qua \(D\).
\end{enumerate}
Bổ đề được chứng minh.

Bây giờ ta chứng minh kết luận 2 bằng cách chia hai trường hợp.

\paragraph{Trường hợp 1: Có hai đỉnh của hình bình hành nằm trên cùng một cạnh của \(P\).}
Giả sử \(A,B\) cùng nằm trên một cạnh. Khi trượt \(AB\) dọc cạnh đó, có thể
tìm hình bình hành \(A'B'CD\) diện tích không đổi sao cho một điểm trong bốn
điểm trở thành đỉnh của \(P\). Do đó kết luận đúng trong trường hợp này.

\paragraph{Trường hợp 2: Không có hai điểm nào cùng nằm trên một cạnh.}
Khi đó \(A,B,C,D\) lần lượt nằm trên bốn cạnh khác nhau \(a,b,c,d\) của
\(P\), và đều không phải đầu mút cạnh. Gọi \(E\) là tâm hình bình hành.

Nếu \(a\parallel c\) và \(b\parallel d\), vị trí của \(E\) được xác định.
Ta có
\[
  \area{ABCD}=4\area{\triangle ABE}=2AE\cdot H(B,AC).
\]
Giữ nguyên đường chéo \(AC\), điều chỉnh \(BD\): qua \(B\) kẻ đường \(e\)
song song \(AC\). Một đầu mút \(V\) của cạnh \(b\) chắc chắn nằm trên \(e\)
hoặc ở phía của \(e\) xa \(E\). Dời \(B\) về phía \(V\) một lượng vi mô thành
\(B'\), đồng thời dời \(D\) theo hướng ngược lại thành \(D'\). Rõ ràng
\[
  H(B',AC)\ge H(B,AC),
\]
nên \(\area{AB'CD'}\) không nhỏ hơn \(\area{ABCD}\). Cứ dời vi mô cùng hướng
cho đến khi \(B\) hoặc \(D\) đến một đỉnh của \(P\), kết luận được chứng minh.

Trường hợp tổng quát hơn là \(a\not\parallel c\) hoặc \(b\not\parallel d\).
Không mất tính tổng quát giả sử \(b\not\parallel d\). Gọi giao điểm hai đường
thẳng chứa \(b,d\) là \(F\). Gọi hai đầu mút của \(b\) là \(B_1,B_2\), trong
đó \(B_1\) nằm giữa \(F\) và \(B_2\); hai đầu mút của \(d\) là \(D_1,D_2\),
trong đó \(D_1\) nằm giữa \(F\) và \(D_2\). Vùng tam giác tạo bởi hai tia
\(B_2B_1,D_2D_1\) và đoạn \(B_1D_1\) được gọi là vùng thứ nhất; vùng vô hạn
tạo bởi hai tia \(B_1B_2,D_1D_2\) và đoạn \(B_2D_2\) được gọi là vùng thứ hai.

Vì \(A,B,C,D\) đều nằm trên bao lồi \(P\), trong hai điểm \(A,C\), đúng một
điểm nằm trong vùng thứ nhất và điểm còn lại nằm trong vùng thứ hai. Giả sử
\(C\) nằm trong vùng thứ nhất, \(A\) nằm trong vùng thứ hai. Ta sẽ cố định
\(A\), điều chỉnh \(B,C,D\) để làm tăng diện tích hình bình hành.\footnote{Nếu
cố định \(C\) rồi điều chỉnh \(A,B,D\), lập luận phía sau không còn đúng.}

Gọi \(l\) là đường thẳng chứa cạnh \(b\), \(g\) là đường thẳng chứa cạnh
\(d\). Trên đoạn \(c\), lấy hai điểm \(C'\) và \(C''\) đối xứng nhau qua
\(C\), với \(CC'\) tiến về \(0\). Gọi \(E'\) và \(E''\) lần lượt là trung điểm
của \(AC'\) và \(AC''\). Gọi \(B'D'\) và \(B''D''\) lần lượt là các đường
chia đôi của \((l,g,E')\) và \((l,g,E'')\).

Do tứ giác có hai đường chéo cắt nhau tại trung điểm là hình bình hành,
\(AB'C'D'\) và \(AB''C''D''\) đều là hình bình hành, tâm lần lượt là \(E'\) và
\(E''\). Chỉ cần chứng minh
\[
  \area{AB'C'D'}+\area{AB''C''D''}>2\area{ABCD},
\]
thì trong hai hình bình hành mới, ít nhất một hình có diện tích lớn hơn
\(\area{ABCD}\), và kết luận theo đó được chứng minh.

Vì
\[
  \area{ABCD}=4\area{\triangle ADE},\quad
  \area{AB'C'D'}=4\area{\triangle AD'E'},\quad
  \area{AB''C''D''}=4\area{\triangle AD''E''},
\]
ta chỉ cần chứng minh mệnh đề tương đương
\[
  \area{\triangle AD'E'}+\area{\triangle AD''E''}>
  2\area{\triangle ADE}.
\]

Do \(E'\) và \(E''\) đối xứng nhau qua \(E\), bổ đề 2 cho biết \(D'\) và
\(D''\) đối xứng nhau qua \(D\). Mặt khác, vì \(EE'\parallel g\) và \(g\) cắt
\(c\), nên \(EE'\) cắt \(DD'\). Gọi giao điểm là \(X\). Xét trường hợp thứ tự
\(A,D,X\) là ngược chiều kim đồng hồ; trường hợp thuận chiều kim đồng hồ tương
tự, còn trường hợp \(D\) trùng \(X\) cũng chứng minh được bằng cách tương tự.

Khi đó, phân rã các diện tích thành các tam giác nhỏ theo cấu hình của hình
gốc, ta có
\[
\begin{aligned}
&\area{\triangle AD'E'}+\area{\triangle AD''E''}\\
&=(\area{AD''XE'}-\area{AD'D''}-\area{D'E'E''}-\area{XE''D'})\\
&\quad+(\area{AD''XE'}-\area{AE'E''}-\area{E''D'D''}-\area{XE''D})\\
&=2\area{AD''XE'}-2\area{XE''D'}
  -\area{AD'D''}-\area{AE'E''}-\area{D'E'E''}-\area{E''D'D''},
\end{aligned}
\]
trong khi
\[
\begin{aligned}
2\area{\triangle ADE}
&=2(\area{AD''XE'}-\area{ADD''}-\area{AEE'}-\area{D'EE''}
  -\area{EDD'}-\area{XE''D})\\
&=2\area{AD''XE'}-2\area{XE''D'}
  -\area{AD'D''}-\area{AE'E''}-\area{D'E'E''}-\area{ED'D''}.
\end{aligned}
\]
Lấy hiệu, các hạng tử chính triệt tiêu:
\[
  \area{\triangle AD'E'}+\area{\triangle AD''E''}
  -2\area{\triangle ADE}
  =\area{ED'D''}-\area{E''D'D''}>0.
\]
Do đó
\[
  \area{\triangle AD'E'}+\area{\triangle AD''E''}>2\area{\triangle ADE}.
\]
Suy ra
\[
  \area{\triangle AD'E'}>\area{\triangle ADE}
  \quad\text{hoặc}\quad
  \area{\triangle AD''E''}>\area{\triangle ADE},
\]
tương ứng
\[
  \area{AB'C'D'}>\area{ABCD}
  \quad\text{hoặc}\quad
  \area{AB''C''D''}>\area{ABCD}.
\]
Nói cách khác, nếu \(ABCD\) nằm trên biên \(P\) nhưng không điểm nào là đỉnh
của \(P\), thì \(ABCD\) chắc chắn không phải hình bình hành diện tích lớn nhất
trong \(P\). Kết hợp các trường hợp, ta chứng minh được: có thể chọn một trong
\(A,B,C,D\) là đỉnh của \(P\) mà không làm giảm diện tích.\footnote{Lý do
không thể cố định \(C\) rồi điều chỉnh \(A,B,D\) là khi đó vị trí tương đối
của \(D'',D,D'\) thay đổi; khi ấy \(\area{\triangle AE'D'}+
\area{\triangle AE''D''}\) không luôn lớn hơn
\(2\area{\triangle AED}\).}

\paragraph{Tiểu kết của kết luận 2.}
Kết luận 2 khó hơn kết luận 1 nhiều. Nhìn lại chứng minh, ta phải chia nhiều
trường hợp; may là phần lớn là trường hợp đặc biệt và tương đối dễ chứng minh.
Trường hợp tổng quát \(b\not\parallel d\) là khó nhất.

Trong chứng minh, ta cố định \(A\), chọn hai điểm đối xứng cách \(C\) một
khoảng vô cùng bé ở hai phía, rồi điều chỉnh \(ABCD\) theo hai hướng thành
\(AB'C'D'\) và \(AB''C''D''\). Bổ đề 2 cho một tính chất đơn giản: \(B',B''\)
đối xứng qua \(B\), và \(D',D''\) đối xứng qua \(D\). Nhờ tính chất này, ta
thu được
\[
  \area{AB'C'D'}+\area{AB''C''D''}>2\area{ABCD}.
\]

Đây là một ứng dụng khá phức tạp của phương pháp giới hạn. Lần hiệu chỉnh này
dùng khoảng cách chứ không dùng góc, vì khi hiệu chỉnh theo khoảng cách, thay
đổi của các đại lượng tương đối dễ nắm bắt; nếu dùng góc, các biểu thức lượng
giác có lẽ sẽ phức tạp hơn nhiều. Điều này cho thấy vai trò quan trọng của khả
năng quan sát trong phương pháp giới hạn.

Ngoài ra, như trong ví dụ 2, chứng minh dùng cách hiệu chỉnh hai lần rồi cộng.
Vì sao không hiệu chỉnh một lần? Có hai lý do chính:
\begin{enumerate}
  \item Ta không biết hiệu chỉnh về phía nào sẽ làm hàm mục tiêu tốt hơn; không
  phải hướng nào cũng tốt hơn.
  \item Hiệu giữa ``tổng hai hàm mục tiêu sau hiệu chỉnh'' và ``hai lần hàm
  mục tiêu ban đầu'' rất dễ tính. Khi tách thành các tổng diện tích tam giác,
  phần lớn hạng tử bị triệt tiêu.
\end{enumerate}
Nếu chỉ hiệu chỉnh một lần, ta phải tính lúc nào nên hiệu chỉnh theo hướng nào;
so với hiệu chỉnh hai lần, đó là một công việc phụ thêm, và chứng minh cũng sẽ
phức tạp hơn. Kỹ thuật hiệu chỉnh hai lần rồi cộng là một thủ pháp thường gặp
trong chứng minh bằng phương pháp giới hạn.

\subsection{Tìm cực đại bằng hàm bậc hai}

Tổng kết như vậy có vẻ hơi sớm: ta mới chứng minh được kết luận 2, nhưng các
hình bình hành thỏa kết luận 2 vẫn còn vô hạn. Bài toán ban đầu chưa được giải
trọn vẹn. Có thể ta sẽ lập tức phỏng đoán: nếu yêu cầu hai trong bốn điểm
\(A,B,C,D\) là đỉnh của \(P\), liệu giá trị tối ưu có giữ nguyên không? Thoạt
nhìn có vẻ có thể dùng phương pháp tương tự để chứng minh: theo kết luận 2,
giả sử \(A\) là đỉnh của \(P\) còn \(B,C,D\) nằm trên các cạnh. Vì nếu cố định
\(B,C\) hoặc \(D\), điểm \(A\) có thể bị điều chỉnh ra ngoài \(P\), nên ta chỉ
có thể cố định \(A\) và điều chỉnh \(B,C,D\). Nhưng nếu \(A\) nằm trong vùng
thứ nhất ở lập luận trên, kết luận 2 không còn đúng. Vì thế rất có thể diện
tích lớn nhất của hình bình hành có hai điểm là đỉnh của \(P\) nhỏ hơn diện
tích lớn nhất của hình bình hành nằm trong \(P\).\footnote{Thực tế có phản ví
dụ; phỏng đoán này không phải khó chứng minh mà là sai.} Phỏng đoán thất bại.

Ta chỉ có thể dùng kết luận 2 hiện có để tìm cực đại. May mắn là chỉ cần kiến
thức về hàm bậc hai. Cách làm như sau.

Trước hết để tính toán đơn giản, thực hiện các phép biến đổi sau:
\begin{enumerate}[label=\alph*)]
  \item Quay \(P\) sao cho cạnh \(c\) song song trục \(y\).
  \item Tịnh tiến \(P\) sao cho hoành độ của \(A\) và hoành độ của \(c\) đối
  nhau. Giả sử \(A=(-k,0)\) và \(c\) có phương trình \(x=k\).
  \item Co giãn đồng dạng tâm \((0,0)\) sao cho \(A\) trở thành \((0,-1)\).
\end{enumerate}
Phép quay và tịnh tiến không làm thay đổi giá trị tối ưu. Sau phép co giãn,
diện tích lớn nhất trở thành \(1/k^2\) lần diện tích ban đầu, nên chỉ cần nhân
kết quả tìm được với \(k^2\) để thu được đáp án trước biến đổi.

Sau chuẩn hóa, ta có thể mô tả cấu hình như sau. Đường chứa \(A\) là
\(x=0\), cạnh \(c\) là \(x=1\), và
\[
  A=(0,0)\quad\text{trong hệ tọa độ cục bộ của công thức diện tích}.
\]
Điểm \(E\), tâm hình bình hành, chạy trên trục \(y\); đặt \(E=(0,t)\). Một
cạnh chứa \(B\) có phương trình
\[
  y=k_1x+b_1,
\]
cạnh chứa \(D\) có phương trình
\[
  y=k_2x+b_2.
\]
Khi đó
\[
  B=(0,b_1),\qquad C=(1,2t),
\]
điểm đối xứng với \(B\) qua \(E\) nằm trên đường qua \(C\) song song với cạnh
\(AB\), nên
\[
  F=(1,k_1+2t-b_1),
\]
và hoành độ của \(D\) là
\[
  D.x=\frac{b_1+b_2-2t}{k_1-k_2}.
\]
Theo hình trong bài gốc,
\[
  \area{ABCD}=4\area{\triangle DEC}
  =4(\area{CFDE}-\area{CFD})
  =4(\area{DFE}+\area{CFE}-\area{CFD}).
\]
Ta tính từng phần:
\[
\begin{aligned}
\area{DFE}
&=\frac12 DF\cdot h\\
&=\frac12\left[(1-D.x)\sqrt{k_1^2+1}\right]\,
  \frac{k_1\cdot0-t+(2t-b_1)}{\sqrt{k_1^2+(-1)^2}}\\
&=\frac12\left(1-\frac{b_1+b_2-2t}{k_1-k_2}\right)(b_1-t)\\
&=\frac{(k_1-k_2+2t-b_1-b_2)(b_1-t)}{2(k_1-k_2)}.
\end{aligned}
\]
Tiếp theo,
\[
\begin{aligned}
\area{CFE}-\area{CFD}
&=\frac12 CF(1-0)-\frac12 CF(1-D.x)\\
&=\frac12 CF\cdot D.x\\
&=\frac12\left[2t-(k_1+2t-b_1)\right]\,
  \frac{b_1+b_2-2t}{k_1-k_2}\\
&=\frac{(b_1-k_1)(b_1+b_2-2t)}{2(k_1-k_2)}.
\end{aligned}
\]
Do đó
\[
\begin{aligned}
\area{ABCD}
&=2\cdot
\frac{(2t+k_1-k_2-b_1-b_2)(b_1-t)
 +(b_1-k_1)(b_1+b_2-2t)}
 {k_1-k_2}\\
&=2\cdot
\frac{-2t^2+(5b_1+b_2-3k_1+k_2)t
 +(2b_1k_1-b_1k_2-2b_1b_2+k_1b_2-2b_1^2)}
 {k_1-k_2}\\
&=2\cdot
\frac{-2t^2+(k_1+k_2+b_1+b_2)t-(k_2b_1+k_1b_2)}
 {k_1-k_2}.
\end{aligned}
\]
Đây đúng là một hàm bậc hai theo \(t\). Từ kiến thức cơ bản về hàm bậc hai,
để \(\area{ABCD}\) đạt cực đại, \(t\) hoặc nằm ở một đầu mút của miền xác
định, hoặc bằng hoành độ trục đối xứng:
\[
  t=\frac{k_1+k_2+b_1+b_2}{4}.
\]
Khi \(t\) là đầu mút của miền xác định, một trong \(B,C,D\) là đỉnh của \(P\),
trường hợp này dễ xử lý. Khi
\[
  t=\frac{k_1+k_2+b_1+b_2}{4},
\]
tọa độ của \(C\) là
\[
  C=\left(1,\frac{k_1+k_2+b_1+b_2}{2}\right).
\]
Điểm này đúng là trung điểm của đoạn nối hai giao điểm \(c\cap b\) và
\(c\cap d\). Vì thế trong thuật toán thật, ta không cần thực hiện các phép
quay, tịnh tiến, co giãn phức tạp: chỉ cần tính tọa độ của hai giao điểm
\(c\cap b\) và \(c\cap d\), rồi dùng công thức trung điểm để lấy tọa độ \(C\).

Đến đây bài toán đã được giải trọn vẹn. Vì bài viết này chủ yếu trình bày ứng
dụng của phương pháp giới hạn, tác giả không nêu thuật toán cụ thể; người đọc
có thể tự hoàn thiện dựa trên quá trình phân tích.

\section{Kết luận}

Qua phân tích cẩn thận ba ví dụ trên, hy vọng người đọc đã dần hiểu ý nghĩa và
cách dùng phương pháp giới hạn, đồng thời thấy sức mạnh của nó: ngắn gọn và
thực dụng.\footnote{Sau khi phân tích xong, các thuật toán đều rất gọn.} Có
thể nói phương pháp giới hạn là con đường ngắn để giải các bài toán tối ưu
hình học phẳng.

Trong chứng minh, phương pháp giới hạn đòi hỏi nền tảng hình học phẳng khá
vững, nên sử dụng có độ khó nhất định. Tinh túy của nó cần được tự lĩnh hội
trong quá trình phân tích; muốn nắm linh hoạt lại càng cần tích lũy kinh
nghiệm.

\section*{Lời cảm ơn}

Tác giả chân thành cảm ơn Li Shi và He Lin đã giúp đỡ trong quá trình nghiên
cứu ví dụ 2. Tác giả cũng chân thành cảm ơn thầy Xiang Qizhong cùng các bạn
trong nhóm yêu thích tin học đã đọc bài viết và đưa ra nhiều ý kiến, đề xuất
sửa đổi quý báu.

\section*{Tài liệu tham khảo}

Không có.

\end{document}
